\documentclass[lettersize,journal]{IEEEtran}
\usepackage{amsmath,amsfonts}
\usepackage{algorithmic}
\usepackage{array}
\usepackage[caption=false,font=normalsize,labelfont=sf,textfont=sf]{subfig}
\usepackage{textcomp}
\usepackage{stfloats}
\usepackage{url}
\usepackage{verbatim}
\usepackage{graphicx}
\hyphenation{op-tical net-works semi-conduc-tor IEEE-Xplore}
\def\BibTeX{{\rm B\kern-.05em{\sc i\kern-.025em b}\kern-.08em
    T\kern-.1667em\lower.7ex\hbox{E}\kern-.125emX}}
\usepackage{balance}
\usepackage{cite}
\usepackage{tabularx}
\hyphenpenalty=25
\tolerance=10000
\allowdisplaybreaks[4]
\usepackage{tcolorbox}
\newcommand{\ciao}[1]{{\setlength\fboxrule{0pt}\fbox{\tcbox[colframe=black,colback=white,shrink tight,boxrule=0.5pt,extrude by=1mm]{\small #1}}}}

\setlength{\abovedisplayskip}{1ex}
\setlength{\belowdisplayskip}{1ex}
\setlength{\textfloatsep}{0.75\baselineskip plus 2pt minus 2pt}
\setlength{\dbltextfloatsep}{0.75\baselineskip plus 2pt minus 2pt}

\usepackage[colorlinks=true,linkcolor=black,citecolor=blue,urlcolor=blue,]{hyperref}

\begin{document}

\title{\vspace{-0.5cm}
Multi-Area Frequency Security Region with Dynamic Tie-Line Mutual Assistance
}
\author{Yiping Yuan, Member, \textit{IEEE}, Liqian Zhu, Student Member, \textit{IEEE}, Zhenyuan Zhang, Senior Member, \textit{IEEE}, Weihao Hu, Senior Member, \textit{IEEE}, Zhe Chen, Fellow, \textit{IEEE}
}
\IEEEaftertitletext{\vspace{-2.5\baselineskip}}

\markboth{Journal of \LaTeX\ Class Files,~Vol.~18, No.~9, September~2020}%
{Multi-area frequency security region with dynamic tie-line mutual assistance}

\maketitle

\begin{abstract}
As power systems move toward low-inertia operation, frequency security increasingly depends on where inertia and damping resources are located. Existing single-area $H$-$D$ support sets identify feasible inertia-damping bounds, but they do not capture dynamic assistance through inter-area tie-lines. This letter proposes a multi-area frequency security region (MA-FSR) framework for quantifying regional operating bounds under coupled frequency dynamics. The framework combines a decoupled analytical approximation with a state-space dynamic search that accounts for tie-line capacity and synchronizing stiffness limits. Case studies on the IEEE two-area Kundur system show that dynamic mutual assistance can substantially expand feasible $H$-$D$ regions, especially in the disturbed area. The MA-FSR therefore provides a compact way to evaluate how interconnection strength changes regional frequency-security margins.
\end{abstract}

\begin{IEEEkeywords}
System frequency response, multi-area power systems, tie-line mutual assistance, frequency security region.
\end{IEEEkeywords}

\vspace{-0.25cm}

\section{Introduction}
\IEEEPARstart{T}{he} rapid integration of converter-interfaced generation (CIG) is reducing the rotational inertia available in modern power grids. Frequency-security constraints, including rate of change of frequency (RoCoF) and frequency nadir, have therefore become binding limits in unit commitment, dispatch, and planning studies~\cite{8918440,10168982,markovic2019optimal,9765359}. The single-area $H$-$D$ support set~\cite{manuscript} provides a compact representation of feasible equivalent inertia ($H$) and damping ($D$) parameters. However, practical power systems are organized as interconnected control areas. Their regional frequency responses are reshaped by tie-line power exchange, which cannot be represented by an isolated single-area model.

This letter extends the frequency security region concept to multi-area power systems. During a contingency, the disturbed area can import power from healthy areas, thereby reducing its local frequency excursion. We study this effect using two complementary evaluation methods. Option A is a decoupled analytical approximation that treats tie-line exchange as an equivalent reduction in step disturbance. Option B is a dynamic simulation-based bisection search that preserves the coupled multi-state interaction and the nonlinear tie-line capacity limit. The resulting multi-area frequency security region (MA-FSR) is visualized in the $H$-$D$ plane, and the effects of tie-line capacity ($C_{12}$) and synchronizing stiffness ($T_{12}$) are quantified. The simulation repository is open-sourced at \url{https://github.com/Divas1234/frequencyregions.git}.

\vspace{-0.25cm}
\section{Multi-Area Frequency Dynamics \& Modeling}

Consider a two-area interconnected power system in which Area 1 is disturbed by a step power imbalance $\Delta P_1$, while Area 2 provides dynamic mutual assistance through the tie-line.

\subsection{Option A: Decoupled Approximation}
In the decoupled approach, tie-line exchange is approximated as a static contribution to the area power imbalance. For Area $i$ connected to neighboring areas, the effective disturbance $\Delta P_{eff,i}$ is defined as
\vspace{-0.15cm}
\begin{equation}
\Delta P_{eff,i} = \Delta P_i + \sum_{j \in \Omega_i} \kappa_{ij} C_{ij}
\vspace{-0.15cm}
\end{equation}
where $C_{ij}$ is the tie-line transmission capacity, and $\kappa_{ij} \in [-1, 1]$ represents the direction and magnitude of the assisting power flow. The single-area analytical bounds in~\cite{manuscript} are then evaluated with $\Delta P_{eff,i}$ to construct decoupled $H$-$D$ regions.

\subsection{Option B: Dynamic State-Space Modeling}
To capture the transient interaction, the interconnected system is represented by a five-state model. For each area $i \in \{1, 2\}$, the coupled swing equation is
\vspace{-0.15cm}
\begin{equation}
2 H_i \dot{\Delta f_i} = \Delta P_{m,i} - D_i \Delta f_i - \Delta P_{dist,i} + \Delta P_{tie,i}
\vspace{-0.15cm}
\end{equation}
where $\Delta P_{tie,1} = \Delta P_{tie}$ and $\Delta P_{tie,2} = -\Delta P_{tie}$. The governor-turbine systems are modeled as:
\vspace{-0.15cm}
\begin{subequations}
\begin{align}
\Delta P_{m,i} &= - K_{m,i} \Delta f_i + x_{g,i} \\
\dot{x}_{g,i} &= \frac{1}{T_{g,i}} \left[ - (1/R_i - K_{m,i}) \Delta f_i - x_{g,i} \right]
\end{align}
\end{subequations}
where $T_{g,i}$ is the turbine-governor time constant, $R_i$ is the droop gain, and $K_{m,i}$ is the high-pressure fraction. The tie-line power flow is governed by
\vspace{-0.15cm}
\begin{equation}
\dot{\Delta P}_{tie} = 2 \pi f_0 T_{12} (\Delta f_2 - \Delta f_1)
\vspace{-0.15cm}
\end{equation}
subject to the capacity saturation constraint $|\Delta P_{tie}| \leq C_{12}$, where $T_{12}$ is the synchronizing coefficient and $f_0 = 50$ Hz is the nominal frequency.

Defining the state vector $x = [\Delta f_1, x_{g,1}, \Delta f_2, x_{g,2}, \Delta P_{tie}]^T$ and input vector $u = [\Delta P_{dist,1}, \Delta P_{dist,2}]^T$, the unsaturated dynamics can be written in state-space form as
\begin{equation}
\dot{x} = A x + B u
\end{equation}
where the state matrix $A$ and input matrix $B$ are
{\small
\begin{equation*}
A = \begin{bmatrix}
-\frac{D_1 + K_{m,1}}{2H_1} & \frac{1}{2H_1} & 0 & 0 & \frac{1}{2H_1} \\
-\frac{1/R_1 - K_{m,1}}{T_{g,1}} & -\frac{1}{T_{g,1}} & 0 & 0 & 0 \\
0 & 0 & -\frac{D_2 + K_{m,2}}{2H_2} & \frac{1}{2H_2} & -\frac{1}{2H_2} \\
0 & 0 & -\frac{1/R_2 - K_{m,2}}{T_{g,2}} & -\frac{1}{T_{g,2}} & 0 \\
-2\pi f_0 T_{12} & 0 & 2\pi f_0 T_{12} & 0 & 0
\end{bmatrix}
\end{equation*}
}
{\small
\begin{equation*}
B = \begin{bmatrix}
-\frac{1}{2H_1} & 0 \\
0 & 0 \\
0 & -\frac{1}{2H_2} \\
0 & 0 \\
0 & 0
\end{bmatrix}^T
\end{equation*}
}
When the tie-line power flow reaches the transmission limit, $\Delta P_{tie}$ is clamped to $\pm C_{12}$ and the corresponding state derivative $\dot{\Delta P}_{tie}$ is set to zero.

\vspace{-0.25cm}
\section{Multi-Area Frequency Security Region}
The MA-FSR is defined as the set of regional inertia ($H_i$) and damping ($D_i$) parameters that satisfy both local and system-wide frequency-security constraints:
\vspace{-0.15cm}
\begin{subequations}
\begin{align}
\max |\Delta f_i(t)| &\leq \gamma_{Nadir,i}, \quad \forall i \\
\max |\dot{\Delta f}_i(t)| &\leq \gamma_{RoCoF,i}, \quad \forall i \\
\zeta_i &\leq 1.0, \quad \forall i
\end{align}
\end{subequations}
Because the initial RoCoF occurs at $t = 0^+$, before the tie-line flow deviates, the local RoCoF boundary remains independent of tie-line parameters:
\vspace{-0.15cm}
\begin{equation}
H_{i,RoCoF} = \frac{\Delta P_{dist,i}}{2 \gamma_{RoCoF,i}}
\vspace{-0.15cm}
\end{equation}

Following the single-area $H$-$D$ support-set derivation in~\cite{manuscript}, the dynamic feasibility of Area $i$ can be separated into local stability bounds and frequency-performance bounds. Let $D_i$ be fixed during boundary tracing. The damping-dependent stability interval is written as
\vspace{-0.15cm}
\begin{subequations}
\begin{align}
H_{i,stab}^{-}(D_i) &=
\frac{T_{g,i}}{2}\left(\eta_i-\sqrt{\eta_i^2-(D_i+K_{m,i})^2}\right),\\
H_{i,a}^{+}(D_i) &=
\frac{T_{g,i}}{2}\left(\eta_i+\sqrt{\eta_i^2-(D_i+K_{m,i})^2}\right),\\
H_{i,b}^{+}(D_i) &= \frac{T_{g,i}}{2}(D_i+K_{m,i}),\\
H_{i,stab}^{+}(D_i) &= \min\!\left[H_{i,a}^{+}(D_i),\,H_{i,b}^{+}(D_i)\right],
\end{align}
\vspace{-0.15cm}
\label{eq:stab_bounds}
\end{subequations}
where $\eta_i = D_i-K_{m,i}+2/R_i$. These bounds are inherited from the damping-ratio condition $0 \leq \zeta_i \leq 1$ and prevent the search from accepting frequency trajectories that satisfy nadir limits but violate local modal damping requirements.

The nadir boundary is denoted by $H_{i,Nadir}(D_i)$. Under Option A, it is obtained by evaluating the closed-form single-area nadir expression with $\Delta P_{eff,i}$. Under Option B, it is obtained from the coupled time-domain trajectory. The lower and upper MA-FSR envelopes are therefore
\vspace{-0.15cm}
\begin{subequations}
\begin{align}
H_i^{-}(D_i) &= \max\!\left[
H_{i,RoCoF},\, H_{i,stab}^{-}(D_i),\, H_{i,Nadir}(D_i)
\right],\\
H_i^{+}(D_i) &= H_{i,stab}^{+}(D_i).
\end{align}
\vspace{-0.15cm}
\label{eq:mafsr_envelope}
\end{subequations}
The area is feasible at damping $D_i$ only if $H_i^{-}(D_i) \leq H_i^{+}(D_i)$. This construction turns each regional frequency-security requirement into an explicit geometric boundary in the $H$-$D$ plane.

To compute the MA-FSR boundary under Option B, a bisection search is performed:
\begin{enumerate}
    \item For a given damping value $D_i$ in Area $i$, initialize the regional inertia interval $[H_{min}, H_{max}]$.
    \item At $H_{mid} = \frac{1}{2}(H_{min} + H_{max})$, simulate the five-state model using fourth-order Runge-Kutta integration.
    \item Track $\Delta f_j(t)$ for all areas $j \in \{1, 2\}$ and test the nadir condition $\max |\Delta f_j(t)| \leq \gamma_{Nadir,j}$.
    \item If both areas satisfy the nadir constraints, set $H_{max} = H_{mid}$; otherwise, set $H_{min} = H_{mid}$.
    \item The search terminates when $|H_{max} - H_{min}| < \epsilon$ (where $\epsilon = 10^{-3}$), yielding the critical inertia limit $H_{i,crit}(D_i)$.
\end{enumerate}
A quadratic regression is then fitted:
\vspace{-0.15cm}
\begin{equation}
H_{i,crit}(D_i) = c_0 + c_1 D_i + c_2 D_i^2
\vspace{-0.15cm}
\end{equation}
The fitted $H_{i,crit}(D_i)$ is used as $H_{i,Nadir}(D_i)$ in~\eqref{eq:mafsr_envelope}. The feasible operating region is then obtained by sweeping $D_i$ and retaining the interval between $H_i^{-}(D_i)$ and $H_i^{+}(D_i)$.

\vspace{-0.25cm}
\section{Case Studies and Analysis}
The IEEE two-area Kundur test system is used for verification. Both areas have nominal droop gains $1/R_i = 33.33$ p.u. and turbine parameters $T_{g,i} = 0.25$ s and $K_{m,i} = 0.35$. The nadir and RoCoF thresholds are $\gamma_{Nadir,i} = 0.5$ Hz and $\gamma_{RoCoF,i} = 0.5$ Hz/s. Area 1 is subjected to a local step contingency of $\Delta P_1 = 3.5$ p.u., while Area 2 has $\Delta P_2 = 2.0$ p.u. The nominal tie-line parameters are $T_{12} = 4.0$ p.u. and $C_{12} = 1.0$ p.u.

\begin{figure*}[!t]
  \centering
  \includegraphics[width=\textwidth]{multiarea_comparison_dynamic.pdf}\vspace{-0.1cm}
  \caption{MA-FSR comparison between isolated and interconnected states under Option B.}\vspace{-0.25cm}
  \label{fig:comparison}
\end{figure*}

\subsection{MA-FSR Expansion under Interconnection}
Fig.~\ref{fig:comparison} compares the feasible operating regions of Area 1 and Area 2 in the $H$-$D$ plane. The isolated case uses $C_{12} = 0$, whereas the interconnected case uses $C_{12} = 1.0$ p.u. In Area 1, which experiences the heavier disturbance, interconnection shifts the nadir boundary downward. At nominal damping, the required inertia decreases from approximately $99.6$ s to $59.5$ s, which substantially enlarges the feasible region. In Area 2, the interconnected nadir boundary reaches the search floor of $0.05$ s because Area 2 is supported by Area 1's nominal inertia ($H_1 = 8.0$ s). Its feasible region is therefore governed mainly by the local RoCoF requirement, $H_2 \geq 4.0$ s.

\subsection{Impact of Tie-Line Capacity and Stiffness}
Fig.~\ref{fig:capacity} shows how Area 1's security region changes with tie-line capacity $C_{12}$ and synchronizing stiffness $T_{12}$. As $C_{12}$ increases from 0 to 4.0 p.u. (Fig.~\ref{fig:capacity}a), the critical inertia boundary decreases progressively. Once $C_{12}$ exceeds the peak tie-line power flow of approximately 2.5 p.u., additional capacity no longer expands the feasible region. Increasing $T_{12}$ (Fig.~\ref{fig:capacity}b) accelerates inter-area power transfer, reduces the transient frequency drop in Area 1, and further enlarges the feasible $H$-$D$ region.

\begin{figure*}[!t]
  \centering
  \subfloat[Tie-line capacity sweep]{\includegraphics[width=0.495\textwidth]{multiarea_capacity_impact.pdf}}\hfill
  \subfloat[Stiffness sweep]{\includegraphics[width=0.495\textwidth]{multiarea_stiffness_impact.pdf}}\vspace{-0.1cm}
  \caption{Impact of tie-line capacity $C_{12}$ and stiffness $T_{12}$ on Area 1's feasible boundary.}\vspace{-0.25cm}
  \label{fig:capacity}
\end{figure*}

\vspace{-0.25cm}
\section{Conclusion}
This letter extends the frequency security region framework to multi-area power systems. By combining a decoupled approximation with a nonlinear dynamic simulator, the proposed MA-FSR characterizes regional $H$-$D$ operating bounds under tie-line mutual assistance. The IEEE two-area Kundur case shows that interconnection can substantially relax the disturbed area's inertia requirement, while the benefit saturates when tie-line capacity exceeds the peak assisting power flow. The MA-FSR provides a compact tool for coordinating regional frequency-control parameters and assessing multi-area scheduling security under explicit tie-line limits.

\bibliographystyle{IEEEtran}
\bibliography{IEEEabrv,reference}

\end{document}
